\section{Semantic Context Triangulation}

\subsection{Purpose}

The project can now begin semantic reconstruction in a constrained way. This section connects the structural and abstract phonetic models to corpus metadata: site, region, object type, material, and iconographic symbol labels. The goal is to discover whether a candidate reading unit behaves like a name, title, place, office, commodity, emblem, or administrative formula.

This is not translation. A context association such as ``unit \texttt{176} is often found on bull-associated seals'' does not mean that \texttt{176} reads ``bull''. It means that any future reading of \texttt{176} must explain that context.

\subsection{Method}

The script \texttt{scripts/semantic-context-triangulation.py} uses the 607 medium/high phonetic-readiness reconstructions. For each slot filler and reading unit, it counts contextual distributions across:

\begin{itemize}
  \item region and site;
  \item artifact type;
  \item material;
  \item iconographic symbol label;
  \item semantic frame.
\end{itemize}

For each unit/context pair, the model computes observed count, expected count, lift, log-lift, and an association score. The score rewards repeated support, purity within the unit's occurrences, and lift above corpus baseline. The output prioritizes units already marked as strong onomastic anchors, cross-frame phonetic anchors, or affix/formula anchors.

\subsection{Corpus-Wide Result}

\begin{table}[htbp]
\centering
\begin{tabular}{lr}
\toprule
\textbf{Metric} & \textbf{Value} \\
\midrule
Ready reconstructions tested & 607 \\
Reading units profiled & 627 \\
Strong onomastic units profiled & 7 \\
Usable iconography labels & 482 \\
Unknown iconography labels & 125 \\
Strong context associations & 60 \\
Moderate context associations & 186 \\
Symbol-linked semantic candidates & 33 \\
Semantic candidates exported & 113 \\
\bottomrule
\end{tabular}
\caption{Semantic context triangulation summary.}
\label{tab:semantic-context-summary}
\end{table}

\subsection{Prime-Entity Semantic Leads}

The first strong semantic target remains the \texttt{002 X 740} prime-entity frame. The current best leads are:

\begin{itemize}
  \item \texttt{240-482} (\texttt{A02-R16}): Harappa/tablet-heavy; likely administrative or institutional label.
  \item \texttt{176} (\texttt{R02}): Mohenjo-daro plus Lothal; bull/gaur seal context; possible name, office, lineage, or emblem-linked identity.
  \item \texttt{235-222} (\texttt{M01-R31}): bull seals with Chanhu-daro concentration; emblem-linked name/title candidate.
  \item \texttt{032-220} (\texttt{A07-A09}): bull seals across sites; cross-site emblem/name candidate.
  \item \texttt{590-390} (\texttt{A08-M02}): steatite seals; bull-linked where known; name/title/place candidate.
  \item \texttt{240-176} (\texttt{A02-R02}): steatite seals; sparse but phonetic-ready candidate.
  \item \texttt{235-840-032} (\texttt{M01-R10-A07}): elephant/bull seals; sparse prime-entity candidate.
\end{itemize}

\subsection{Important Non-Onomastic Leads}

Some high-scoring units are not likely to be proper names. They may be titles, suffixes, object labels, or administrative formulae:

\begin{itemize}
  \item \texttt{240-520} strongly clusters with Harappa tablets and bull-labeled contexts. It may be a terminal administrative expression rather than a personal name.
  \item \texttt{240-740-090} is associated with tablet contexts and a \texttt{Cros} symbol label. This makes it a strong object-class or administrative-ending candidate.
  \item \texttt{032} is notably Chanhu-daro-heavy, suggesting a possible regional, scribal, or administrative specialization.
  \item \texttt{140} has a strong goat-labeled subgroup while also appearing in bull contexts, so it should be treated as a classifier/title/formula marker until image review resolves the context.
\end{itemize}

\subsection{Interpretation}

The project has now reached semantic proto-reconstruction. The safest current semantic template for many seal inscriptions is:

\begin{center}
\texttt{CLASSIFIER/TITLE + ENTITY SLOT + TERMINAL/FINAL}
\end{center}

For the prime frame this is:

\begin{center}
\texttt{002 + ENTITY + 740}
\end{center}

The entity may be a personal name, title, lineage, place, institution, commodity, or emblem-linked identity. The context model now lets us choose which possibility to test first for each filler.

\subsection{Next Empirical Test}

The most important next step is image-grounded validation. The corpus symbol labels are useful but not enough. For each strong anchor, we need to inspect the actual seal or inscription images and ask:

\begin{itemize}
  \item Does the same unit consistently occur with the same iconographic motif?
  \item Does it remain stable across site, material, and object type?
  \item Do minimal alternations preserve the same iconographic and structural context?
  \item Does a proposed semantic reading predict all occurrences, including the awkward ones?
\end{itemize}

The immediate priority set is \texttt{176}, \texttt{235-222}, \texttt{032-220}, \texttt{590-390}, and \texttt{240-482}. These give us the first real bridge from structure toward meaning.

\subsection{Outputs}

\begin{itemize}
  \item \texttt{outputs/semantic\_context\_summary.csv}
  \item \texttt{outputs/anchor\_context\_profiles.csv}
  \item \texttt{outputs/anchor\_context\_correlations.csv}
  \item \texttt{outputs/semantic\_reconstruction\_candidates.csv}
  \item \texttt{outputs/iconography\_semantic\_matrix.csv}
  \item \texttt{outputs/semantic\_context\_triangulation.tex}
\end{itemize}
